\subsection*{MT-S505. Real-axis Artin holomorphy in the specified $S_5$ family}

Retain the extension and local hypotheses of MT-S501 and the
notation of MT-S504. Conditional on $M_H$, every irreducible
Artin factor of this extension is holomorphic at every real
point of $(0,1)$, including the centre. In particular
\[
 \operatorname{ord}_{1/2}\zeta_H
  \ge \operatorname{ord}_{1/2}\zeta_F
                  +\operatorname{ord}_{1/2}L(s,\eta).
 \tag{S505.1}
\]
The difference in (S505.1) is a nonnegative even integer.
If it is zero, every irreducible factor is automorphic
and satisfies Artin holomorphy and GRH, by MT-S504.
Thus no additional central nonvanishing hypothesis is
needed for the real-axis holomorphy assertion itself.

\begin{proof}
Write $m_5=\operatorname{ord}_{1/2}L(s,\rho_5)$,
allowing a negative order for a meromorphic pole.
If $m_5<0$, the sufficient condition $m_5\le0$
of MT-S504 holds. That theorem proves automorphy
and forces $m_5=0$, a contradiction. Hence $m_5\ge0$.
Meromorphicity then means that $L(s,\rho_5)$
is holomorphic at $s=1/2$. This argument does
not assume its central value is nonzero.

For completeness the other rank-five factor
is also holomorphic at the centre for an
unconditional elementary reason. Put
$J=\operatorname{Gal}(L/F)$, of order $20$.
Its Sylow $5$-subgroup $P$ is normal, since
its number divides $4$ and is $1$ modulo $5$.
The normalizer of $P$ in $S_5$ has order $20$;
thus $J$ is that normalizer. In the standard
affine action on $\mathbb F_5$, its elements
have cycle types
\[
       1\text{ once},\quad 5A\text{ four times},\quad
       2B\text{ five times},\quad 4A\text{ ten times}.
\]
Using Calegari's character table, arXiv:1112.1152v1,
p.~3, the average over these elements is $1$
for the trivial character and for $\rho_5\otimes\eta$,
and $0$ for $\eta,\rho_4,\rho_4\otimes\eta,\rho_5,\rho_6$.
For example the two rank-five averages are
\[
 \frac{5+4\cdot0+5\cdot1+10\cdot(-1)}{20}=0,
 \qquad
 \frac{5+4\cdot0+5\cdot1+10\cdot1}{20}=1.
\]
Frobenius reciprocity therefore proves the
permutation identity and its all-prime
Artin factorization
\[
 \operatorname{Ind}_{J}^{S_5}\mathbf1
          =\mathbf1\oplus(\rho_5\otimes\eta),\qquad
 \zeta_F(s)=\zeta_{\mathbb Q}(s)L(s,\rho_5\otimes\eta).
 \tag{S505.2}
\]
For $0<s<1$, the alternating series
\[
 A(s)=\sum_{n\ge1}(-1)^{n-1}n^{-s}
     =\sum_{j\ge1}\bigl((2j-1)^{-s}-(2j)^{-s}\bigr)
\]
converges and is strictly positive. Analytic continuation
of the elementary identity from $s>1$ gives
$A(s)=(1-2^{1-s})\zeta_{\mathbb Q}(s)$ on this interval.
The first factor is negative, so
$\zeta_{\mathbb Q}(s)<0$ throughout $(0,1)$.
Thus (S505.2) expresses the twisted rank-five
factor as a holomorphic quotient on the whole
interval, without any GRH assumption.

The factors of dimensions $4$ and $6$ are
entire by the cited Calegari Theorem~1.1,
and the linear factors are holomorphic
on $(0,1)$. Finally for $\rho_5$ at a
real point other than $1/2$, signed Brauer
induction and $M_H$ give a product and
quotient of nonzero holomorphic Hecke values,
so it is holomorphic there too.
This proves the real-axis assertion for
every irreducible factor.

Taking central orders in (S504.1) gives
(S505.1) with difference $m_5$. The real
realizability/root-number-$+1$ and
meromorphic functional-equation parity
argument already checked in MT-S504
makes $m_5$ even. The equality case is
exactly the sufficient condition of
MT-S504. All conclusions follow.
\end{proof}

This settles the real-axis portion of the
Artin pole gate for this precise $S_5$
family. Possible critical-line poles away
from the real axis are not ruled out in
the remaining $m_5\ge2$ case. Geometrically,
(S505.2) is the degree-six permutation
Artin motive of the finite cover, while
(S505.1) is an effectivity statement for
the central coefficient of the relative
rank-five divisor. It concerns that
coefficient, not the effectivity of the
whole Artin divisor required by
Q-ARTIN-POS and the cote~1/SGA5 trace programme.
